\section{Proof Sketches and Technical Notes}
\label{app:proofs}

\subsection{Need-sufficiency architecture
(Section~\ref{sec:need_sufficiency})}

\paragraph{Proposition~\ref{prop:below_suff_saturation}
(Below-Sufficiency Saturation).}
Below-sufficiency benchmark for dimension $d_k$ is
$\min(\mathrm{actual}_k / s_k, 1)$, bounded above by $1$
regardless of the number of capabilities reporting
$\mathrm{actual}_k \geq s_k$.  This proposition is deliberately
narrow: it is the full extent of the need-sufficiency
architecture's structural anti-stuffing contribution.  General
M5 capability-stuffing is identified in
Section~\ref{sec:capability_stuffing} as an instance of
structural avoidance; the resolution direction is a proof-burden
taxonomy with claim-and-challenge back-pressure, deferred to
follow-up work (Open Question~\ref{oq:capability_stuffing}).

\paragraph{Proposition~\ref{prop:s1_sufficiency}
(S1--Sufficiency Connection).}
Below-sufficiency sacrifice events fail S1's non-degrading
alternative condition.  Key step: retaining $X$ (not making the
trade) leaves the agent below functionality on dimension $d_k$.
S1 requires that retaining $X_n$ was a non-degrading
alternative.  When retaining $X$ means falling below sufficiency
on $d_k$, S1's condition fails directly.

\subsection{Privacy-minimality
(Section~\ref{sec:sacrifice_model})}

\paragraph{Theorem~\ref{thm:privacy} strategy.}
The five properties (P1--P5) are proved independently:
P1~(no interior access): by construction of the observation
tuple.
P2~(consent): by the summation structure of the lower bound.
P3~(discretization): by bounding channel bandwidth with trade
frequency.
P4~(commitment composability): by composition
with~\citep[Definition~6]{lasser2026exo}.
P5~(third-party observability): by the public nature of market
transactions.

\subsection{Aggregate lower bound
(Section~\ref{sec:lower_bound})}

\paragraph{Theorem~\ref{thm:aggregate_bound} strategy.}

\emph{Part~A (bundle-level, S0--S2, poset-independent bundles):}
\begin{enumerate}
\item Per-event bound from calibration chain (S0 + S1 + S2):
  $\volR(Y_n) \geq \Ui[_{n}](Y_n) \geq \Ui[_{n}](X_n) \geq
  \Delta\volL(X_n)$.
\item Poset-independence of bundles
  (Definition~\ref{def:category_partition}) gives M4 additivity:
  $\volR(\bigcup Y_n) = \sum_n \volR(Y_n) \geq \sum_n
  \Delta\volL(X_n)$.
\end{enumerate}

\emph{Part~B (per-capability, additionally S3--S5):}
\begin{enumerate}
\setcounter{enumi}{2}
\item S5 gives directly $\alpha_{n,c} \leq \volR(c)/\volR(Y_n)$
  (a share inequality, not a decomposition identity), yielding
  $\volR(c) \geq \alpha_{n,c} \cdot \Delta\volL(X_n)$ without
  requiring within-bundle poset-independence.
  S3 (bundle coherence) ensures $\volR(Y_n) > 0$ so the division
  is well-defined.
\item Per-capability max-attribution across overlapping events
  (Proposition~\ref{prop:monotone}).
\item Summation over capabilities in $U$.
\end{enumerate}

Key technical challenges:
\begin{itemize}
\item \emph{Poset-independence (Part~A):} the category partition
  $\mathcal{C}$ (Definition~\ref{def:category_partition}) must be
  a genuine poset partition.  Set-disjointness is necessary but
  not sufficient (Remark~\ref{rem:poset_indep}).
\item \emph{Decomposition validity (S5, Part~B):} hedonic
  regression~\citep{rosen1974hedonic} produces attribution weights,
  not $\volR$ share lower bounds.  S5 is the explicit bridge.
\item \emph{Calibration (S0):} the weakest assumption in the chain
  but the hardest to verify empirically (Open
  Question~\ref{oq:calibration}).
\end{itemize}

\subsection{Axiom inheritance
(Section~\ref{sec:axiom_inheritance})}

\paragraph{Theorem~\ref{thm:axiom_inheritance} strategy.}
Axiom-by-axiom analysis:
M1--M5 are direct from properties of summation/max over
non-negative terms.
M6 failure~(a), $\volR$ itself: constructive counterexample
showing that merging groups changes exercise indicators
(alternative pathways make previously load-bearing capabilities
dormant).
M6 failure~(b), $\volRlower$: constructive counterexample showing
that merging groups with overlapping bundle categories changes
per-capability attributions.
Corollary~\ref{cor:not_self_balancing} follows from~(a), not~(b).

\subsection{Connection to observational individuation}

The transfer argument is direct.  Active agents generate the OI
floor
of~\citep[Corollary~2.1]{lasser2026horizon} as a byproduct of
participation.  The OI floor contributes to $\volR$ independently
of the sacrifice channel.  For agents who are both active and
trading, the $\volRcomb$ diagnostic
(Appendix~\ref{app:exercise_indicator}) includes both the OI
floor and the sacrifice lower bound.
